\documentclass[twocolumn]{aastex631}
\usepackage{amsmath}
\usepackage{booktabs}
\shorttitle{Selection-aware SDSS BPT/sSFR study}
\shortauthors{NebulaMind}
\begin{document}

\title{Broad Optical BPT Galaxies and Catalog Specific Star Formation in SDSS DR17: A Selection-Aware Pilot Matched-Control Study}
\author{NebulaMind Research Autopilot}
\affiliation{Public SDSS DR17 data only}

\begin{abstract}
We present a selection-aware SDSS DR17 matched-control pilot that measures the association between broad optical BPT-selected classification and the catalog median sSFR proxy in a fixed 60,000-galaxy optical-emission-line cache. Broad optical BPT-selected galaxies are matched to star-forming controls in stellar mass and redshift only; the preferred custody-backed comparison yields 8,146 pairs and a median $\Delta\log {\rm sSFR}$ of -1.309 dex, with a bootstrap 95\% confidence interval of [-1.334,-1.283] dex. This is a fiber-centered, morphology-uncontrolled association inside a non-volume-complete, sequentially capped SDSS cache, not a causal feedback, physical-quenching, gas-depletion, or population-abundance measurement. The candidate custody file inventories the retained cache, pilot JSON, and matched-pair CSV that support the matched-offset result; public selection-cascade counts are retained only as context unless accompanied by separate query receipts. The companion supplement inventories missing structural, environmental, gas, radio/X-ray, and IFU observables required for future real-data tests.
\end{abstract}

\keywords{galaxies: active --- galaxies: star formation --- galaxies: evolution --- surveys --- methods: statistical}

\section{Question and claim boundary}
This paper addresses a narrow association-only question within a low-redshift SDSS DR17 optical emission-line denominator: do broad optical BPT-selected galaxies have lower catalog median sSFR proxy than mass--redshift matched star-forming controls? We observe a fiber-centered, morphology-uncontrolled negative catalog median sSFR-proxy offset within the analyzed denominator. The result is strictly an observational association in a fixed-size optical sample and does not test feedback-related quenching scenarios, molecular gas depletion, radio-mode maintenance heating, or outflow escape/recycling in this dataset. It is not a causal result. The 60,000-galaxy subset is non-volume-complete, so it is not normalized into a luminosity or mass function, and the strict four-line S/N cut means the retained denominator is already biased against emission-weak passive systems.


The present scope also excludes morphology or aperture controls, structural-proxy matching, Seyfert/LINER separation, bolometric accretion-luminosity proxies, gas-mass measurements, environment labels, and time-domain or duty-cycle modelling. BPT line ratios classify optical excitation, not black-hole accretion power; as seen in previous literature, retired stellar populations ionized by hot post-AGB stars, as well as low-ionization nuclear emission-line region (LINER)-like ionization and extended low-ionization emission-line regions, can contaminate broad low-ionization classes and mimic active-nucleus signatures \citep{cidfernandes2011,stasinska2008,stasinska2015,belfiore2016}. For that reason the paper uses the phrase ``broad optical BPT-selected galaxies'' and treats stricter Seyfert/LINER separation as future work rather than as an interchangeable label. Because structural proxies were not retained in the 60,000-galaxy cache, the present optical denominator cannot separate the measured offset from bulge-fraction, concentration-index, \texttt{fracDeV}, or central-velocity-dispersion associations. Because the sample is restricted to $0.02<z<0.12$, the standard local BPT demarcations are used here without any redshift-evolution correction.

\subsection{Scope and limitations}
The association reported here is defined inside a capped, selection-limited optical denominator. It is not a volume-complete census, and it does not include morphology, aperture fraction, group membership, halo mass, gas mass, or bolometric accretion-luminosity proxies as matching variables. Those missing dimensions are relevant follow-up requirements, but they are not part of the present inference. In particular, the absent structural proxies mean the observed fiber-centered offset remains inseparable from bulge-related or aperture-related associations in the retained cache.

\section{Missing observables for future causal inference}
The present SDSS-only analysis is deliberately restricted to an optical association pilot, so the observables needed for any causal interpretation are not measured here. The remaining requirements are morphology and structural proxies, aperture-fraction control, group or halo membership, CO/HI gas masses, radio and X-ray proxies, resolved IFU kinematics, and matched simulation comparisons passed through the same selection function. These are the real-data inputs that the companion supplement inventories as future follow-up targets; they are not results of this paper.

\section{Data and shared selection}
The data backbone is public SDSS DR17 spectroscopy, photometry, emission-line measurements, and catalog physical-property estimates \citep{york2000,sdssdr17,brinchmann2004}. The pilot analysis sample is a fixed 60,000-galaxy subset selected sequentially by \texttt{specObjID}. It is a selection-limited pilot subset used to estimate the association within the available SDSS cache, not a volume-limited census. Because \texttt{specObjID} ordering follows SDSS targeting and plate/MJD bookkeeping, this non-random subset is not a random sky sample and introduces survey-plate and sky-coverage bias. The strict public four-line S/N$\geq3$ eligible parent contains 249,917 galaxies, so the pilot sample covers 24.0\% of that strict parent. Because the subset is fixed and non-volume-limited, it cannot be used to derive absolute volume densities, luminosity functions, or any population-normalized abundance.
Over the redshift interval $0.02<z<0.12$, the SDSS 3-arcsec fiber subtends roughly 1.2--6.5 kpc, so the catalog median sSFR proxy comparison is fiber-centered rather than global.
Because the 3-arcsec fiber samples only the central regions at low redshift, the catalog-derived total sSFR proxy is an aperture-extrapolated quantity; the fixed 3-arcsec aperture systematically misses extended star-forming disks at low redshift \citep{kewley2005}. If broad optical BPT hosts are more bulge-dominated than the star-forming controls, the central fiber measurement can inflate the observed offset relative to a galaxy-wide star-formation comparison.
The stellar-mass and sSFR values are taken from the public MPA-JHU-style value-added table \texttt{galSpecExtra}, using its catalog median estimators \texttt{lgm\_tot\_p50} and \texttt{specsfr\_tot\_p50} after joining \texttt{SpecObj}, \texttt{galSpecInfo}, and \texttt{PhotoObj}. Although \texttt{PhotoObj} was joined in the catalog backbone, structural quantities such as \(R_{90}/R_{50}\), \texttt{fracDeV}, \texttt{petroR50}, and \texttt{petroR90} were not retained in the 60,000-galaxy cache, so morphology cannot be controlled in this cycle. Those are low-redshift SDSS catalog estimates, not rederived line-by-line physical measurements \citep{brinchmann2004,sdssdr17,york2000}. We use variance-normalized Euclidean matching, with each coordinate standardized by its sample standard deviation before distance calculation, because the feature space is only two variables, $(\log M_\star,z)$, so the rule stays transparent and the resulting nearest-neighbor control remains easy to interpret as an association baseline.

\begin{deluxetable*}{lrrr}
\tabletypesize{\scriptsize}
\tablecaption{Context selection cascade for the flagship analysis sample.\label{tab:selection}}
\tablehead{\colhead{Selection stage} & \colhead{Public DR17 rows} & \colhead{Cached rows} & \colhead{Retention vs. spectro-z parent}}
\startdata
SpecObj GALAXY, 0.02<z<0.12 & 501,060 & -- & 100.0\% \\
plus galSpecInfo/PhotoObj/galSpecExtra and mass/sSFR bounds & 416,554 & -- & 83.1\% \\
plus galSpecLine join & 416,554 & -- & 83.1\% \\
four BPT lines with valid flux measurements (\texttt{ivar} $> 0$) & 373,445 & 60,000 & 74.5\% \\
four BPT lines S/N>=3 & 249,917 & 60,000 & 49.9\% \\
four BPT lines S/N>=5 & 176,523 & 42,446 & 35.2\% \\
four BPT lines S/N>=10 & 91,768 & 22,311 & 18.3\% \\
\enddata
\tablecomments{Counts are read-only public SDSS DR17 count-query context plus the cached analysis table. Retention is shown as a percentage of the spectro-z parent. Cached rows are shown only where the cache applies. The 416,554-to-373,445 drop when requiring \texttt{ivar} $> 0$ reflects the removal of rows with unusable line-flux uncertainties; this table does not distinguish masking, edge-of-chip loss, or missing spectral coverage. The sharp retention drop at higher S/N mainly reflects preferential loss of passive galaxies from the emission-line denominator. The present custody receipt inventories the retained cache and matched-analysis files, not the public query logs behind every parent-count row, so non-cached parent rows are contextual selection diagnostics rather than independent retained results.}
\end{deluxetable*}

The selection is not neutral with respect to star formation. The strict four-line S/N cut preferentially removes emission-weak passive systems, while the sequential \texttt{specObjID} cap can inherit SDSS plate/MJD and sky-coverage structure. The present candidate custody file inventories the retained cache and matched analysis artifacts, but not the public SQL count logs behind every selection-cascade context number. Those context counts should therefore be accompanied by an explicit query receipt before they are promoted to independent measured results in a journal submission. This caveat does not alter the inventoried 60,000-row cache or the 8,146-pair matched-offset result.

\section{Classification and matching}
BPT classes are computed from H$\alpha$, H$\beta$, [O~III]$\lambda5007$, and [N~II]$\lambda6584$ using standard demarcations \citep{baldwin1981,kewley2001,kauffmann2003bpt,kewley2006} (see Figure~\ref{fig:bpt}). The custody-backed analysis denominator contains 39,553 star-forming galaxies, 12,234 intermediate/composite galaxies, 8,146 broad optical BPT-selected targets, and 67 unclassified objects. The intermediate/composite galaxies are retained in denominator counts but are not part of the star-forming control pool used for matching. The 67 unclassified objects are retained in denominator counts for completeness but excluded from control pairing; they do not enter the 8,146-pair estimate because matching is performed only for broad optical BPT-selected targets with valid star-forming controls. Here, the star-forming control pool is defined as objects below the Kauffmann et al.\ (2003) demarcation, which is a conservative optical cut and may still include weak or obscured active-nucleus contaminants. Each broad optical BPT-selected galaxy is matched to the nearest star-forming control by variance-normalized Euclidean distance in standardized $(\log M_\star,z)$ space, with replacement. In the preferred estimate, this yields 100\% target coverage (8,146 of 8,146 targets matched), and the unrestricted Euclidean match has median absolute separations of 0.0045 dex in $\log M_\star$ and 0.00021 in redshift, so the association still inherits any mismatch in structure or fiber coverage between the two populations. No maximum mass--redshift caliper is imposed because the retained pilot estimate prioritizes complete coverage of all 8,146 custody-backed targets; imposing a caliper would create an additional selection function whose excluded targets are not separately inventoried in this candidate package. Matching is not performed in morphology, aperture fraction, halo mass, gas mass, accretion-luminosity proxy, or duty-cycle phase; these missing dimensions define follow-up requirements.
Here, ``broad optical BPT-selected'' means the inclusive optical-emission-line class under the standard BPT demarcations. Stricter Seyfert/LINER separation is a required follow-up analysis rather than a retained result in this provenance-limited manuscript.

\begin{figure*}
\centering
\includegraphics[width=0.72\textwidth]{../figures/fig-bpt.pdf}
\caption{BPT line-ratio diagram for the SDSS DR17 analysis denominator. The matched controls are paired in stellar mass and redshift only, not in morphology, so the diagram verifies the optical-excitation classes used for matching but does not by itself establish an accretion-based interpretation.}
\label{fig:bpt}
\end{figure*}

\section{Matched-control result}
The preferred broad optical BPT comparison gives a large negative catalog median sSFR-proxy offset for the broad optical BPT-selected galaxies relative to star-forming controls.
    \par\noindent\textbf{Morphology and aperture caveat.} Without controlling for structural morphology or aperture fraction, a median $\Delta\log {\rm sSFR}$ (target minus matched control) of -1.309 dex is observed within this fixed-size, morphology-uncontrolled optical denominator and fiber-centered matched comparison. Because the spectroscopy samples only the central 3-arcsec region (1.2--6.5 kpc here) and the match does not control morphology, structural proxies, or aperture fraction, the observed fiber-centered sSFR offset is highly degenerate with the known correlation between stellar mass and galaxy morphology and the transition from disk-dominated to bulge-dominated systems, including bulge-to-total ratio, bulge prominence, or central velocity dispersion associations \citep{schawinski2010,bluck2014,piotrowska2022}. The lack of concentration-index or \texttt{fracDeV}-style structural matching limits the result's ability to separate bulge-linked structural associations from excitation-linked associations. Single-fiber measurements can miss substantial extended star-forming disks, so spatially resolved integral-field spectroscopy is required to resolve the aperture-morphology degeneracy \citep{penny2018,cheung2016,bundy2015,canodiaz2016}. Within the present data, the association is therefore still not separable from a morphology or bulge-fraction association. The robustness interval in Table~\ref{tab:robust} is a 95\% confidence interval on the median fiber-centered offset.

\begin{deluxetable*}{lrrrr}
\tabletypesize{\scriptsize}
\tablecaption{Provenance-retained matched catalog-sSFR offset.\label{tab:robust}}
\tablehead{\colhead{Variant} & \colhead{$N$ pairs} & \colhead{Median $\Delta\log {\rm sSFR}$} & \colhead{95\% interval} & \colhead{Interpretation}}
\startdata
Broad optical BPT-selected targets, S/N$\geq3$, nearest SF control with replacement & 8,146 & -1.309 & [-1.334,-1.283] & Preferred association estimate \\
\enddata
\tablecomments{$\Delta\log {\rm sSFR}$ is target minus matched star-forming control. This row is retained because it traces to \texttt{SDSS\_AGN\_SFR\_PILOT\_20260708T122000Z/analysis\_results.json} and \texttt{matched\_agn\_sf\_pairs.csv}, both inventoried in \texttt{provenance/REAL\_DATA\_SOURCE\_CUSTODY.json}. Caliper, no-replacement, stricter-S/N, and Seyfert-like sensitivity values are not repeated here because the present custody file does not inventory a separate result artifact for those variants. All values are conditional on the optical emission-line denominator.}
\end{deluxetable*}

\begin{figure*}
\centering
\includegraphics[width=0.86\textwidth]{../figures/fig-matched-offsets.pdf}
\caption{Distribution of matched-pair catalog-sSFR offsets for the preferred broad optical BPT-selected galaxy minus nearest star-forming control estimate ($N=8{,}146$ pairs, without a maximum mass--redshift caliper). The preferred estimate is strong within this denominator, but the present custody file does not inventory separate artifacts for caliper, no-replacement, stricter-S/N, or Seyfert-like sensitivity variants.}
\label{fig:offsets}
\end{figure*}

\section{Interpretation}
The result is directly measured in the capped sample and remains falsifiable within the stated denominator. The matched-offset distribution is shown in Figure~\ref{fig:offsets}. Because the comparison is fiber-centered and selection-limited, the fiber-centered offset remains a denominator-level association statement rather than a galaxy-wide causal inference; the 3-arcsec fiber can preferentially sample central bulge light over extended star-forming disks. Matching on mass and redshift alone leaves morphology uncontrolled, and the fixed 3-arcsec fiber can under-sample extended disk star formation at low redshift; this is a known source of central-to-global mismatch \citep[e.g.,][]{kewley2005,harrison2017,ellison2021}. The retained cache also lacks direct structural quantities, even though SDSS concentration index, \texttt{fracDeV}-style decompositions, and bulge/disk catalogues are standard routes for separating morphology from excitation class in follow-up work \citep{strateva2001,mendel2014,bluck2014,piotrowska2022}. The most robust conclusion is therefore: broad optical BPT classification is associated with lower catalog median sSFR proxy in this fixed-size, selection-limited, morphology-uncontrolled 60,000-galaxy pilot sample, not a causal result. Any mechanistic interpretation requires additional real data, including morphology and aperture controls, Seyfert/LINER separation, bolometric accretion-luminosity proxy, gas mass, environment, and time-domain/duty-cycle modelling; see Supplement Sections 5.1 and 5.7 for the neighbor-rank/fiber-collision and CO/HI requirements. These are missing observables in the present catalog and are required for future mechanism tests.
The choice of variance-normalized Euclidean matching is deliberate: with only two standardized coordinates, it preserves a simple nearest-neighbor control rule without introducing an additional model layer that would not be better constrained by the available data.

\section{Conclusion}
RP-1 is a selection-aware pilot association paper. This analysis is bounded by a fixed 60,000-galaxy subset selected sequentially by \texttt{specObjID}, and lacks morphological, structural, and aperture-fraction controls. Without controlling for structural morphology or aperture fraction, the reported fiber-centered -1.309 dex sSFR offset is an association-only measurement within this fixed-size, morphology-uncontrolled optical denominator and is currently indistinguishable from a morphology, bulge-fraction, or aperture-sampling association. Its provenance-retained result is the preferred 8,146-pair, -1.309 dex offset with bootstrap 95\% interval [-1.334,-1.283] dex. The accompanying \emph{Supplementary SDSS Denominator and Proxy Atlas for Galaxy-Evolution Follow-up} is a denominator/proxy atlas and follow-up target list for the missing-observable requirements of future real-data tests. See the supplement's neighbor-rank/fiber-collision and CO/HI entries, summarized in its atlas overview, for the clearest examples of the remaining constraints. No measured result in this paper should be read as a gas-mass, environment-density, or feedback-efficiency estimate.
Future physical validation requires integration with the kinds of measurements used in radio-mode and X-ray maintenance-heating studies \citep{best2005,fabian2012,mcnamara2007,heckmanbest2014,lamassa2013}, molecular and neutral gas studies \citep{xcoldgass2017,xgass2018}, outflow and kinematic studies \citep{veilleux2005,cicone2014,carniani2017,fiore2017}, and simulation comparisons passed through the same selection functions \citep{simba2019,tng2019,eagle2015}, together with the environment/context references \citep{peng2010,ellison2011,piotrowska2022,wetzel2013,dekel2006}; these references are cited as examples of missing observables for future follow-up, not as validation of any mechanism in this SDSS-only denominator. These are missing observables in the present catalog and are required for future mechanism tests. The result remains association-only until morphology, aperture fraction, and the missing multiwavelength or IFU observables are added. The fixed 60,000-galaxy cache remains a non-random selection-limited subset, so the reported offsets should continue to be read as denominator-bound associations rather than population-wide trends within this morphology-uncontrolled optical denominator.

\section*{Data Availability}
This paper uses public SDSS DR17 spectroscopy, photometry, emission-line measurements, and MPA-JHU-style value-added catalog tables only. No proprietary data were used. The 60,000-row cache is derived from the public catalog joins and selection thresholds described above, and the manuscript conclusions remain conditional on the optical-emission-line denominator. No mock, synthetic, fake, placeholder, or toy data were used.

\facilities{SDSS}

\begin{thebibliography}{}
\bibitem[Abdurro'uf et al.(2022)]{sdssdr17} Abdurro'uf, Accetta, K., Aerts, C., et al. 2022, ApJS, 259, 35
\bibitem[Baldwin et al.(1981)]{baldwin1981} Baldwin, J.~A., Phillips, M.~M., \& Terlevich, R. 1981, PASP, 93, 5
\bibitem[Best et al.(2005)]{best2005} Best, P.~N., Kauffmann, G., Heckman, T.~M., et al. 2005, MNRAS, 362, 25
\bibitem[Belfiore et al.(2016)]{belfiore2016} Belfiore, A., Maiolino, R., Maraston, C., et al. 2016, MNRAS, 461, 3111
\bibitem[Bluck et al.(2014)]{bluck2014} Bluck, A.~F.~L., Bruce, V.~A., Pilkington, K., et al. 2014, MNRAS, 441, 599
\bibitem[Brinchmann et al.(2004)]{brinchmann2004} Brinchmann, J., Charlot, S., White, S.~D.~M., et al. 2004, MNRAS, 351, 1151
\bibitem[Bundy et al.(2015)]{bundy2015} Bundy, K., Law, D.~R., Yan, R., et al. 2015, ApJ, 798, 7; ADS bibcode: 2015ApJ...798....7B
\bibitem[Cano-D\'{\i}az et al.(2016)]{canodiaz2016} Cano-D\'{\i}az, M., Maiolino, R., Marconi, A., et al. 2016, ApJL, 818, L14
\bibitem[Cid Fernandes et al.(2011)]{cidfernandes2011} Cid Fernandes, R., Stasi{\'n}ska, G., Schlickmann, S., et al. 2011, MNRAS, 413, 1687
\bibitem[Cheung et al.(2016)]{cheung2016} Cheung, E., Bundy, K., Cappellari, M., et al. 2016, Nature, 533, 504
\bibitem[Ellison et al.(2011)]{ellison2011} Ellison, S.~L., Patton, D.~R., Mendel, J.~T., et al. 2011, MNRAS, 418, 2043
\bibitem[Ellison et al.(2021)]{ellison2021} Ellison, S.~L., Lin, L., Rosario, D.~J., et al. 2021, MNRAS, 501, 4777
\bibitem[Harrison(2017)]{harrison2017} Harrison, C.~M. 2017, Nature Astronomy, 1, 0165
\bibitem[Carniani et al.(2017)]{carniani2017} Carniani, S., Marconi, A., Maiolino, R., et al. 2017, A\&A, 605, A42
\bibitem[Catinella et al.(2018)]{xgass2018} Catinella, B., Saintonge, A., Janowiecki, S., et al. 2018, MNRAS, 476, 875; ADS bibcode: 2018MNRAS.476..875C
\bibitem[Cicone et al.(2014)]{cicone2014} Cicone, C., Maiolino, R., Sturm, E., et al. 2014, A\&A, 562, A21
\bibitem[Dav{\'e} et al.(2019)]{simba2019} Dav{\'e}, R., Angl{\'e}s-Alc{\'a}zar, D., Narayanan, D., et al. 2019, MNRAS, 486, 2827
\bibitem[Dekel \& Birnboim(2006)]{dekel2006} Dekel, A., \& Birnboim, Y. 2006, MNRAS, 368, 2
\bibitem[Fabian(2012)]{fabian2012} Fabian, A.~C. 2012, ARA\&A, 50, 455
\bibitem[Fiore et al.(2017)]{fiore2017} Fiore, F., Feruglio, C., Shankar, F., et al. 2017, A\&A, 601, A143
\bibitem[Heckman \& Best(2014)]{heckmanbest2014} Heckman, T.~M., \& Best, P.~N. 2014, ARA\&A, 52, 589
\bibitem[Kauffmann et al.(2003a)]{kauffmann2003bpt} Kauffmann, G., Heckman, T.~M., Tremonti, C., et al. 2003a, MNRAS, 346, 1055; ADS bibcode: 2003MNRAS.346.1055K
\bibitem[Kewley et al.(2001)]{kewley2001} Kewley, L.~J., Dopita, M.~A., Sutherland, R.~S., Heisler, C.~A., \& Trevena, J. 2001, ApJ, 556, 121
\bibitem[Kewley et al.(2005)]{kewley2005} Kewley, L.~J., Jansen, R.~A., \& Geller, M.~J. 2005, PASP, 117, 227
\bibitem[Kewley et al.(2006)]{kewley2006} Kewley, L.~J., Groves, B., Kauffmann, G., \& Heckman, T. 2006, MNRAS, 372, 961
\bibitem[LaMassa et al.(2013)]{lamassa2013} LaMassa, S.~M., Heckman, T.~M., Ptak, A., \& Urry, C.~M. 2013, ApJL, 765, L33
\bibitem[McNamara \& Nulsen(2007)]{mcnamara2007} McNamara, B.~R., \& Nulsen, P.~E.~J. 2007, ARA\&A, 45, 117
\bibitem[Nelson et al.(2019)]{tng2019} Nelson, D., Springel, V., Pillepich, A., et al. 2019, Computational Astrophysics and Cosmology, 6, 2
\bibitem[Penny et al.(2018)]{penny2018} Penny, S.~J., Davies, R.~L., Houghton, R.~C.~W., et al. 2018, MNRAS, 476, 979
\bibitem[Peng et al.(2010)]{peng2010} Peng, Y.-j., Lilly, S.~J., Kovac, K., et al. 2010, ApJ, 721, 193
\bibitem[Piotrowska et al.(2022)]{piotrowska2022} Piotrowska, J.~M., Bluck, A.~F.~L., Maiolino, R., \& Peng, Y.-j. 2022, MNRAS, 512, 1052
\bibitem[Schawinski et al.(2010)]{schawinski2010} Schawinski, K., Evans, D.~A., Virani, S., et al. 2010, ApJ, 711, 284
\bibitem[Saintonge et al.(2017)]{xcoldgass2017} Saintonge, A., Catinella, B., Tacconi, L.~J., et al. 2017, ApJS, 233, 22; ADS bibcode: 2017ApJS..233...22S
\bibitem[Schaye et al.(2015)]{eagle2015} Schaye, J., Crain, R.~A., Bower, R.~G., et al. 2015, MNRAS, 446, 521
\bibitem[Strateva et al.(2001)]{strateva2001} Strateva, I., Ivezi{\'c}, {\v Z}., Knapp, G.~R., et al. 2001, AJ, 122, 1861
\bibitem[Mendel et al.(2014)]{mendel2014} Mendel, J.~T., Simard, L., Palmer, M., Ellison, S.~L., \& Patton, D.~R. 2014, ApJS, 210, 3
\bibitem[Stasi{\'n}ska et al.(2008)]{stasinska2008} Stasi{\'n}ska, G., Asari, N.~V., Cid Fernandes, R., et al. 2008, MNRAS, 391, L29
\bibitem[Stasi{\'n}ska et al.(2015)]{stasinska2015} Stasi{\'n}ska, G., Costa Duarte, M.~V., Vale Asari, N., Cid Fernandes, R., \& Sodr{\'e}, L. 2015, MNRAS, 449, 559
\bibitem[Veilleux et al.(2005)]{veilleux2005} Veilleux, S., Cecil, G., \& Bland-Hawthorn, J. 2005, ARA\&A, 43, 769
\bibitem[Wetzel et al.(2013)]{wetzel2013} Wetzel, A.~R., Tinker, J.~L., Conroy, C., \& van den Bosch, F.~C. 2013, MNRAS, 432, 336
\bibitem[York et al.(2000)]{york2000} York, D.~G., Adelman, J., Anderson, J.~E., Jr., et al. 2000, AJ, 120, 1579
\end{thebibliography}

\end{document}
